\begin{abstract}
Heterogeneous cloud computing environments require efficient task scheduling to handle volatile workloads, cut operational costs, and reduce carbon footprints. Static heuristics like HEFT and Min-Min work well under specific conditions, but they fail when workload dynamics change or system faults occur. This dissertation presents a dynamic Multi-Armed Bandit (MAB) hyper-heuristic framework for cloud task scheduling. The architecture links a Python controller to a CloudSim Plus simulation backend using standard I/O streaming pipes. The controller manages a low-level heuristic pool of nine algorithms, including classical list schedulers, meta-heuristics, and local search methods. 

We evaluate four MAB selection mechanisms: $\epsilon$-Greedy, Standard UCB1, Discounted UCB, and Scaled Thompson Sampling. The framework uses a workload fingerprinting method based on the Coefficient of Variation ($\text{CoV}$) alongside a batch-relative multi-objective reward function balancing makespan, execution cost, carbon emissions, and resource utilization. We test the system using HPC2N and NASA real-world trace datasets. Experimental results show that Discounted UCB adapts best to non-stationary environments. When subjected to a 90\% capacity drop, Discounted UCB decayed legacy performance memory and shifted policy from meta-heuristics to PEFT within 10 rounds, whereas Standard UCB remained stuck in stale Q-values. Overall, the hyper-heuristic framework achieved up to a 34.3\% reduction in both operational costs and carbon emissions compared to traditional baselines.
\end{abstract}
